Una vareta metàl·lica es desplaça a una velocitat constant $v = 6\ \mathrm{m/s}$ sobre una forca conductora dins un camp magnètic uniforme, $\vec{B} = 0,25\ \mathrm{T}$, perpendicular al pla i en sentit sortint:

\figura[0.55]{fig1.png}

Si suposem que la resistència de la vareta és de $30\ \Omega$ i que la de la forca és negligible, calculeu:

\begin{apartat}{1}
La força electromotriu del corrent induït en el circuit i expliqueu raonadament el sentit de la circulació del corrent.
\end{apartat}

\begin{apartat}{1}
La intensitat del corrent que circula pel circuit i la força que cal fer sobre la vareta, en mòdul, direcció i sentit, per a mantenir la velocitat constant sobre la forca.
\end{apartat}

\nota{Llei d'Ohm, $I = V/R$.}
